\chapter{Percurso Metodológico}

\section{Natureza, abordagem e estratégia da pesquisa}

A pesquisa será de natureza aplicada, com abordagem predominantemente qualitativa e apoio de dados quantitativos descritivos. Será desenvolvida como intervenção pedagógica de grupo único, com produções individuais em diferentes momentos da sequência. Esse desenho permite examinar mudanças, permanências e dificuldades nas explicações dos estudantes, sem pretender isolar um efeito próprio da plataforma frente à mediação docente, às atividades e às interações em sala.

A estratégia principal será uma intervenção pedagógica. \citeonline{damiani2013discutindo} caracterizam esse tipo de investigação pelo planejamento e pela implementação de mudanças ou inovações pedagógicas, acompanhados da análise de seus efeitos. Essa formulação é compatível com o projeto porque os recursos da intervenção serão planejados para uma situação de ensino, aplicados no campo e examinados a partir das evidências produzidas.

A intervenção ganha sentido como uma unidade integrada que articula a sequência didática e a plataforma aos instrumentos de produção de dados e à mediação docente. Separar artificialmente esses elementos seria inadequado, pois, nesta pesquisa, a plataforma adquire o sentido previsto em conjunto com a pergunta proposta, a previsão solicitada e a discussão posterior.

Reconhecer essa unidade não autoriza atribuir qualquer mudança observada a um único componente. As produções dos estudantes entram como evidências construídas em um percurso situado, no qual a infraestrutura, o tempo de aula, o acesso individual à plataforma e as retomadas conceituais moldam a forma como cada explicação aparece. A posição do professor-pesquisador exigirá registro claro das condições de aplicação e prudência na interpretação dos dados.

A definição metodológica final deverá ser discutida com o orientador antes da submissão do projeto e do início da produção de dados. Essa decisão dependerá também das condições confirmadas do campo e das exigências éticas e institucionais aplicáveis.

\section{Contexto e participantes}

O campo previsto é uma turma de Ensino Médio de escola pública na qual o pesquisador atue como professor. A possibilidade inicial é uma turma vinculada à disciplina Projeto Permanente por Afinidade, em escola de tempo integral. Esse cenário não está garantido para o próximo ano letivo, pois a lotação docente e a composição das turmas poderão mudar.

Se necessário, a sequência poderá ser aplicada em outra turma do Ensino Médio sob docência do pesquisador, inclusive em um componente de Matemática, desde que sejam preservadas as condições pedagógicas e éticas definidas para a pesquisa.

Os participantes serão estudantes iniciantes em \textit{machine learning}, sem exigência de programação prévia. A sequência pressupõe apenas que saibam ler tabelas e gráficos elementares, localizar pontos em eixos cartesianos e comparar valores, além de mobilizar noções iniciais de função e probabilidade, ou de distância e média, conforme forem retomadas nos módulos.

A definição da escola e da série, junto ao número de participantes no componente curricular e ao calendário, ocorrerá após confirmar o campo.

\section{Etapas da intervenção}

O primeiro momento corresponderá ao estudo das referências para aperfeiçoar os exemplos sintéticos da plataforma e os instrumentos previstos. Em seguida, a sequência didática será preparada de acordo com a turma e as condições de infraestrutura disponíveis.

Prevê-se que os estudantes trabalhem individualmente, com um computador do laboratório ou, conforme a disponibilidade confirmada, um Chromebook. No campo inicialmente previsto, o laboratório dispõe de 20 máquinas; em outro campo, será necessário verificar as condições disponíveis. O acesso individual permite que cada estudante formule e teste suas próprias hipóteses diante do mesmo experimento. As previsões e explicações de interesse serão anotadas individualmente antes e depois da exploração em três encontros, enquanto os registros do professor indicarão de que modo o acesso individual, a infraestrutura ou as retomadas conceituais ajudam a dimensionar essas produções. Os oito encontros asseguram uma duração viável, exigindo atenção constante ao ritmo da turma.

Ao final de cada encontro, o professor-pesquisador organizará uma síntese breve que recupere os elementos discutidos pela turma.

As sínteses de encerramento não terão a função de produzir uma resposta correta única para ser copiada pelos estudantes. Elas servirão para retomar o que foi observado e clarear relações confusas antes da passagem ao encontro seguinte. Caso a turma confunda uma probabilidade com garantia de ocorrência, por exemplo, a retomada poderá recuperar o papel do limiar e distinguir estimativa de certeza antes de avançar para outro modelo.

O registro docente indicará que a retomada ocorreu e quais dificuldades a motivaram, sem substituir as respostas individuais como fonte de evidência da compreensão. Essa separação ajuda a preservar o vínculo entre a mediação necessária ao ensino e a análise cuidadosa do que cada participante conseguiu elaborar.

A plataforma não registrará respostas ou dados pessoais. As respostas usadas na pesquisa serão produzidas nos instrumentos descritos a seguir.

\section{Instrumentos de produção de dados}

Antes da sequência, os participantes realizarão um diagnóstico individual breve. O instrumento trará um núcleo de quatro ou cinco questões para examinar desde os dados e a finalidade da tarefa até a interpretação de representações, o efeito de alterações e os limites de decisões automatizadas. Ele servirá para registrar o ponto de partida das explicações e não comporá a avaliação escolar.

Nos encontros 2, 4 e 6, cada participante preencherá um registro curto com duas solicitações: indicar o que prevê antes de alterar o controle e explicar o resultado depois da exploração. A escolha desses encontros permite acompanhar uma tarefa de previsão numérica, uma tarefa de classificação e uma tarefa de agrupamento, sem exigir produção escrita em todos os módulos.

O instrumento final individual terá dois blocos. O primeiro retomará o núcleo do diagnóstico por meio de questões equivalentes em operação cognitiva, extensão e dificuldade. Os contextos poderão mudar, mas cada par de questões deverá solicitar a mesma relação conceitual. O segundo bloco será um desafio de transferência. Nele, o estudante analisará uma situação inédita para reconhecer a tarefa e os dados necessários, indicando um modelo compatível e explicando o efeito esperado caso um dado, parâmetro ou regra sofra alterações.

O desafio de transferência não será tratado como equivalente ao diagnóstico. Sua função será examinar a mobilização dos conceitos em outro contexto. Um questionário breve de percepção registrará a utilidade atribuída às atividades e à plataforma, bem como a clareza e o nível de dificuldade encontrados. Essas respostas serão analisadas como apreciação do produto, e não como demonstração de aprendizagem.

Após cada aula, o professor-pesquisador estruturará anotações em tópicos cobrindo as intervenções realizadas e as dificuldades recorrentes, além do tempo empregado e das condições de uso observadas. As perguntas inseridas no site terão função formativa imediata. Como suas respostas não permanecem salvas, não serão usadas diretamente como dados da pesquisa.

Antes da aplicação, uma matriz associará cada questão comparável ao objetivo de aprendizagem correspondente e à atividade desenvolvida, ligando a dimensão da rubrica ao item equivalente no outro instrumento. Os instrumentos serão submetidos à leitura crítica do orientador e de um professor com experiência no nível de ensino. Quando for viável, a clareza dos enunciados será verificada com pessoa de perfil próximo ao público que não participe da pesquisa.

Os instrumentos também precisarão ser compatíveis com o tempo escolar disponível e com o fato de que os participantes serão estudantes do Ensino Médio em uma atividade inserida no cotidiano da turma. Não se pretende aplicar teste extenso que interrompa a sequência ou transforme a intervenção em rotina de avaliação formal.

Perguntas abertas curtas, acompanhadas de gráfico, tabela ou situação-problema, podem oferecer indícios mais úteis que grande quantidade de itens objetivos, desde que o enunciado deixe claro o que se pede. Antes da aplicação, o professor-pesquisador verificará se cada questão está ligada a um objetivo de aprendizagem e a uma atividade efetivamente realizada.

Caso algum conceito não tenha sido abordado nas condições previstas, ele não deverá ser tratado como ausência de compreensão no instrumento final. Essa precaução mantém os instrumentos vinculados ao percurso de ensino e protege a análise contra interpretações que ultrapassem as evidências disponíveis.

\section{Procedimentos de análise}

As produções serão organizadas por códigos, sem identificação nominal dos participantes. A rubrica terá três níveis, variando de 0 a 2, estruturada em quatro dimensões que vão dos dados e da tarefa até a representação conceitual e os limites da decisão. O nível 0 reunirá respostas ausentes, incompatíveis ou sem relação justificável com a situação. O nível 1 indicará relação parcial entre os elementos. O nível 2 indicará uma explicação consistente com o nível introdutório definido para a sequência. Cada classificação será acompanhada por uma justificativa curta.

A primeira etapa consistirá na leitura integral das respostas e na revisão dos exemplos que delimitam os três níveis. Em seguida, estruturar-se-á uma matriz que cruza cada participante e dimensão aos momentos da coleta. Os códigos apoiarão a organização do corpus, sem virar nota para a turma.

Uma resposta poderá apenas indicar que ``a classificação mudou'', enquanto outra relacionará a mudança ao limiar escolhido ou à proximidade entre exemplos. A diferença entre esses modos de explicar será mais informativa que uma oposição simples entre certo e errado.

As questões comparáveis do diagnóstico e do instrumento final serão colocadas em relação. A análise acompanhará se houve manutenção do nível inicial ou mudança para cima e para baixo, sempre lendo a justificativa escrita. Uma tabela de transição por dimensão mostrará quantos participantes seguiram cada percurso, enquanto os registros intermediários situarão o surgimento dessas relações nas aulas. O desafio de transferência será analisado à parte nas dimensões aplicáveis. Frequências absolutas e relativas virão com seus denominadores. Casos de respostas em branco ou de participantes sem par completo entre os dois instrumentos serão informados sem imputação de valores. Essas medidas descreverão o conjunto e não sustentarão inferências estatísticas ou generalizações para outras turmas.

O texto final trará sínteses por dimensão e frequências descritivas acompanhadas de tabelas de transição, além de breves excertos anonimizados que exemplifiquem formas distintas de explicação. Os registros docentes contextualizarão as produções escritas sempre que dúvidas recorrentes e retomadas matemáticas, ou mesmo condições de infraestrutura, interferirem na atividade. O questionário de percepção será apresentado em seção própria e não será combinado com os níveis da rubrica.

A interpretação considerará as condições em que a atividade ocorreu, mas não usará registros do professor para completar uma resposta individual. Quando for viável, um segundo avaliador classificará uma parcela das produções. As divergências serão discutidas com base nas definições e nos exemplos da rubrica, e os ajustes de critério serão registrados antes da análise final.

\section{Aspectos éticos}

A coleta de dados só terá início após a liberação institucional exigida pelos comitês de ética aplicáveis ao projeto. A Resolução CNS nº 510/2016 dispõe sobre pesquisas em Ciências Humanas e Sociais que utilizam dados obtidos diretamente de participantes ou informações identificáveis, enquanto a Resolução CNS nº 466/2012 estabelece diretrizes gerais para pesquisas envolvendo seres humanos \cite{brasil2016cns510,brasil2012cns466}.

Como a intervenção poderá envolver adolescentes, os procedimentos éticos deverão ser definidos conforme a orientação institucional e a tramitação exigida antes da aplicação.

As respostas receberão códigos de anonimização, mantendo eventuais documentos com dados identificáveis em arquivos separados dos instrumentos destinados à análise. Dados brutos não serão inseridos no código da plataforma nem publicados em repositório aberto. A participação será voluntária, garantindo aos estudantes que optarem por não integrar a pesquisa o acesso a atividades pedagógicas equivalentes sem nenhum prejuízo escolar.
